\documentclass[11pt]{article}
\usepackage{amsmath,amssymb}
\usepackage[hidelinks]{hyperref}
\title{What a live timing amendment must inherit from earlier handoffs}
\date{}
\begin{document}
\maketitle

A shared port has two possible pasts. Both allow exactly the same
continuations under the original timing law, and both say that either
scope can transfer now. After the same authorized timing amendment, the
pasts require opposite first transfers to meet a service deadline. Even
an exact old continuation language can be an inadequate interface for a
later amendment. The physical age of an earlier receipt, not an online
search heuristic, makes the difference.

\section{Scope and timing model}
Use the clocked custody port of \#390, where a port service is an
irreversible transfer and the unique mint of an authenticated receipt.
Signers alternate per scope, one effect uses one slot, and old receipts
may remain acceptable through their declared lifetime unless the port
waits for expiry. Future mechanical minimum intervals can change by an
authorized amendment. This amendment cannot make an already minted old
receipt younger or expire it sooner. The paper does not add another
observation supplier or prove that the clock is distributed. The shared
port may keep a faithful local ledger of actual transfers. A policy
operating only on an advertised old certificate does not have that ledger
for free.

\section{A complete fixed-law deadline interface}
At cut zero consider one scope with a fixed future order of signed tags
and a repeat of length \(K\). Let \(b\) be the time of the last actual
physical transfer before the cut, \(E_j\) the first safe time for the
\(j\)-th upcoming reuse constrained by an old mint, for \(0\leq j<K\),
and \(d_j>0\) the prescribed mechanical gap before future effect \(j\).
An unoccupied old tag imposes no bound. After these \(K\) uses, the
ordinary future reuse checks refer only to future transfers. Only slots
\(0,\ldots,h-1\) are available. Define
\[
 e_{-1}=b,\qquad e_j=\max(0,E_j,e_{j-1}+d_j),\qquad p_j=\min(h,e_j).
\]
A proposed trace is lawful iff its first \(K\) existing times satisfy
\(t_j\geq p_j\), its successive gaps meet the \(d_j\), and it obeys the
future reuse rule when that rule applies. Indeed all actual traces
dominate the earliest prefix by induction. Conversely the profile
bounds recover the old expiry at each relevant slot and the first
mechanical gap from the actual previous transfer. Clipping does not
change any test for times below \(h\). Four Lean theorems check the
arithmetic behind this equivalence; the future reuse law is unchanged
in both arms.

For these trace-membership queries, the profile is exact and necessary.
If two profiles differ, take their first differing index, choose the
history with the earlier value and propose its earliest prefix through
that index. It is a feasible witness before \(h\) and cannot occur in
the later history. This conclusion supposes that exact trace membership
is consumed, or contexts can actually test those schedules. A restricted
client concerned only with one aggregate target can have a coarser
interface. Tuples of profiles suffice for a fixed declared menu of
timing laws, and their necessity uses separate tests for each law in
that menu. Extending the menu changes which histories can be identified.
An enlarged horizon can also expose different old expiration times
collapsed by the earlier cap. This is not a strategy-synthesis theorem
for a law that can change unpredictably during a scheduled prefix.

\section{A live counterexample and its repair}
Before cut zero, use mechanical gaps \((1,5)\) for A and \((5,2)\) for B,
indexed by next signer, with reuse waits \(W_A=6\) and \(W_B=7\).
Consider actual histories in slots \(-9,\ldots,-1\):
\[
 \texttt{A.....B.A},\qquad \texttt{A.....AB.}.
\]
Both end at next signer 0 for A and 1 for B. All future schedules
up to six slots have identical legality under the old law: each
scope has local profile \((0,5)\), and they have identical port
conflict and future reuse rules. The uncapped event checker
independently enumerates the same 53 legal shared schedules for both.
The old law does not meet the harder service request below; preserving
that old legal language is the specific interface being tested.

A new authorized law reduces the second impending mechanical gap
on both scopes from five to four. The first gaps do not change.
Two actual transfers on each scope are now required in slots 0--5.
The first transfers must occupy slots 0 and 1: anything later puts its
second at slot six or later. Both scopes are currently ready in either
history. In the first past, A's next use by signer 1 cannot occur
before slot five because its old receipt was minted at \(-1\); B's
signer 0 can repeat at slot four. In the second past, B's signer 0
cannot occur before five because its old mint was at \(-2\);
A's signer 1 can repeat at four. The ports therefore have unique
useful schedules \texttt{BA..BA} and \texttt{AB..AB}, respectively.
Starting the wrong scope or idling loses at least one required
transfer. No controller that observes only the equal old profile,
the equal current readiness and the common amendment can guarantee
the new target, even with randomization under a sure-service law.

For this pair, one bit from a faithful actual-mint record restores the
guarantee: did A's old signer-1 mint still prevent reuse at slot four?
The amended prefix certificate then directs the correct first effect.
It matters that this source is already at the actual port and that
the evidence is available before deciding. The bit is not obtained
from the readiness output. A direct dispatcher with access to the same
ledger gets the same plans. A detector with no received source record
does not acquire it merely because the mechanics changed.

\section{Status and comparison}
\href{https://pmc.ncbi.nlm.nih.gov/articles/PMC9468080/}{LiveLTL}
distinguishes retained execution obligations from the new specification
after an update. The concrete question here is whether an implementation
that has already quotiented old physical evidence can reconstruct the
different admission language of an authorized update. The example and
profile give an exact finite answer for this trusted-clock model;
they do not displace a matched update synthesizer with all actual
receipt times. For Contextual Composition, the issue is the primitive
test family of the profile: preservation at the old test family does
not give reflection when the receiver adds new mechanical contexts.
This is partial R2, not independent transfer or contract acceptance.
\end{document}
